\begin{proof} 1. \\
  まず以下を証明する.\\
  \begin{equation}
    \InjMap{f}{X}{Y} \rightarrow \BijMap{f}{X}{f[X]}
  \end{equation}
  \begin{itemize}
    \item[]\hspace{-0.5cm} ($\InjMap{f}{X}{f[X]}$) \\
      任意に$x_1, x_2 \in X$をとり,
      $<x_1, y'> \in f \land <x_2, y'> \in f$な$y' \in f[X]$の存在を仮定する.\\
      そんな$y'$を一つとり固定.\\
      このとき,
      $x_1, x_2 \in A$,
      $y' \in f[X] \subseteq Y$,
      $<x_1, y'> \in f \land <x_2, y'> \in f$である.\\
      これらを$\InjMap{f}{X}{Y}$に適用すると$x_1 = x_2$.\\
      よって$\InjMap{f}{X}{f[X]}$.
    \item[]\hspace{-0.5cm} ($\SurMap{f}{X}{f[X]}$) \\
      任意に$y \in f[X]$をとる.\\
      このとき,
      $\exists x \in X (f(x) = y)$は真よりそんな$x$を一つとり固定.\\
      そんな$x$の存在より命題は真.
  \end{itemize}
  $f$が単射であると仮定し,
  1.1から1.3を示す.\\
  \begin{itemize}
    \item[]\hspace{-0.5cm} (1.1) \\
      $Y$が可算なので,
      $Y$は有限集合あるいは可算無限集合.\\
      
    \begin{itemize}
     \item[]\hspace{-0.5cm} (有限集合のとき) \\
       $\BijMap{g}{Y}{n}$をみたす自然数$n$が存在するのでそんな$n$をひとつとり固定.\\
       また,
       $\InjMap {g \after f}{X}{n}$より,
       (7.1)を使うと$\BijMap {g \after f}{X}{(g\after f)[X]}$\\
       $h:= g \after f$とおく.\\
       このとき$h[X] \subseteq n$とCor7.1.6.1より,
       $\BijMap{r}{h[X]}{m}$をみたす自然数$m$が存在するので、そんな$m$を一つとり固定.\\
       このとき$\BijMap{r \after h}{X}{m}$が存在する.
       ゆえに$X$は有限集合である.
     \item[]\hspace{-0.5cm} (可算無限集合のとき) \\
       $\BijMap{g}{Y}{\nat}$をみたす$g$が存在する.\\
       また,
       $\InjMap{g \after f}{X}{\nat}$より,
       $h := g \after f$とおくと$\BijMap {h}{X}{h[X]}$.\\
       このとき,
       $h[X] \subseteq \nat$である.\\
       $h[X] \subseteq n$のときは有限集合のときに証明済.\\
       $h[X] \subseteq n$をみたす自然数$n$が存在しないとき,
       Prop7.1.5.2より$\BijMap{r}{h[X]}{\nat}$が存在する.\\
       このとき$\BijMap{s}{X}{\nat}$が存在する.
       ゆえに$X$は可算無限集合.
    \end{itemize}
    \item[]\hspace{-0.5cm} (1.2) \\
      1.1より$Y$が有限集合のとき$X$も有限集合.\\
      1.1より$g \after f[X] \subseteq n$とProp7.1.5.1より,
      $\Card{(g \after f)} = m = \Card{X} \leq n = \Card{Y}.$\\
    \item[]\hspace{-0.5cm} (1.3) \\
      $\Card{X} = \Card{Y} = n$とすると,
      $\BijMap{g}{X}{n}$,$\BijMap{h}{n}{Y}$が存在する.\\
      それらを合成すると,
      $\BijMap{h \after g}{X}{Y}$が得られる.\\
  \end{itemize}
\end{proof}


\begin{proof} 2. \\
  $f$が全射であると仮定し,
  1.1から1.3を示す.\\
  \begin{itemize}
    \item[]\hspace{-0.5cm} (1.1) \\
      $X$が可算なので,
      $X$は有限集合あるいは可算無限集合.\\
    \begin{itemize}
     \item[]\hspace{-0.5cm} (有限集合のとき) \\
       $\BijMap{g}{n}{X}$をみたす自然数$n$が存在するのでそんな$n$をひとつとり固定.\\
       このとき,
       $\SurMap{f \after g}{n}{Y}$.また$h:=f \after g$とする.\\
       任意に$y \in Y$をとる.\\
       $h$は全射より,
       $\exists a \in n(h(a) = y)$は真なので,
       $h^{-1}[\{y\}]$は空でない.\\
       また,
       $h^{-1}[\{y\}] \subseteq n$より,
       $h^{-1}[\{y\}]$の最小元$l$が存在する.\\
       $r:=<y, l>$とおくと,
       $h \after r = id_y$より,
       $\InjMap{r}{Y}{n}$\\
       $n$は可算であり,
       1.(1.1)より$Y$も可算である.
     \item[]\hspace{-0.5cm} (可算無限集合のとき) \\
       有限集合の場合と同様にして,
       $\nat$も可算なので,
       1.(1.1)より$Y$も可算である.
    \end{itemize}
    \item[]\hspace{-0.5cm} (1.2) \\
      1.1より,
      $\InjMap{g \after r}{Y}{X}$.\\
      1(1.2)より,
      $Y$は有限集合で$\Card{Y} \leq \Card{X}$.
    \item[]\hspace{-0.5cm} (1.2) \\
      1.(1.3)と同様にして,
      $\BijMap{h \after g}{Y}{X}$が得られる.\\
  \end{itemize}
\end{proof}
